% Source: https://artofproblemsolving.com/community/c3h1584400_parabolas
% Stable identity: aops-topic-1584400-q01
The parabola is symmetric about the vertical line $x=1$ through its vertex, so the two $x$-intercepts are equally far from it, and their midpoint is on the axis:
\[
\dfrac{a+b}{2}=1 .
\]
Since $a<0$, $|a|=-a$, so $b=-3a$. Then $\dfrac{a-3a}{2}=-a=1$, so $a=-1$ and $b=3$.

\begin{asy}
usepackage("amsmath");
size(6cm);
defaultpen(fontsize(10pt));
import graph;
real f(real x){return -1.25*(x-1)^2+5;}
draw((-2.3,0)--(4.4,0),gray,Arrow(TeXHead));draw((0,-1.6)--(0,6),gray,Arrow(TeXHead));
label("$x$",(4.4,0),E,gray);label("$y$",(0,6),N,gray);
draw(graph(f,-1.35,3.35),blue+linewidth(1));
draw((1,-1.3)--(1,5.6),dashed);
dot((1,5));dot((-1,0));dot((3,0));
label("$(1,5)$",(1,5),NE);label("$(a,0)$",(-1,0),SW);label("$(b,0)$",(3,0),SE);
label("$x=1$",(1,-1.3),S);
label("\mbox{}",(min(currentpicture,user=true)+max(currentpicture,user=true))/2,black); // keeps the text after the figure black when LaTeX typesets the labels inline
\end{asy}

(The parabola is $y=-\dfrac{5}{4}(x-1)^2+5$, which indeed vanishes at $x=-1$ and $x=3$.) Hence
\[
ab=\boxed{-3}.
\]
